%! Designing Experiments — Unit 3, Lesson 3.4 — Slides (Beamer)
\documentclass[aspectratio=169, 11pt]{beamer}
\usepackage{apstats-beamer}   % provides \navyheader{} and \sectionlabel[color]{}

\begin{document}

% TITLE SLIDE (CourseName is not defined in beamer, so the name is literal)
{
\setbeamercolor{background canvas}{bg=navy}
\begin{frame}[plain]
  \vspace{2.8cm}\hspace{0.55cm}%
  \begin{minipage}{0.9\textwidth}
    {\color{goldacc}\bfseries\small LESSON 3.4}\\[0.5cm]
    {\color{white}\bfseries\LARGE Designing Experiments}\\[1.2cm]
    {\color{skymid}\small Statistics~$\cdot$~Unit 3}
  \end{minipage}
\end{frame}
}

% GRADUAL-RELEASE deck: targets → warm-up → I DO divider → one frame per instruction row of the
% notes (Experiment; Block It?; Matched Pairs, which carries the crux) → WE DO (The Zinc
% Lozenge) → spoken debrief → close & START THE HOMEWORK, which is the individual practice.
% There is no group activity and no independent practice set in the notes.
\newcommand{\redink}[1]{{\color{redacc}\bfseries #1}}

% ── LEARNING TARGETS ─────────────────────────────────────────────────────────
\begin{frame}[t, plain]
  \navyheader{Today's targets}
  \sectionlabel{Lesson 3.4}
  By the end of today, I can\ldots
  \begin{itemize}\setlength{\itemsep}{5pt}
    \item name the \textbf{experimental units}, \textbf{factors}, \textbf{treatments}, and
          \textbf{response variable} of an experiment.
    \item explain \textbf{comparison}, \textbf{random assignment}, \textbf{replication}, and
          \textbf{control} --- with a \textbf{control group}, a \textbf{placebo}, and
          \textbf{blinding}.
    \item choose a \textbf{completely randomized}, \textbf{randomized block}, or
          \textbf{matched pairs} design, and say why chance still assigns inside every block.
  \end{itemize}
  \begin{block}{How today works}
    Warm-up $\to$ \textbf{notes together} $\to$ \textbf{one we work as a class} $\to$
    debrief $\to$ \textbf{you start the homework here, alone}.
  \end{block}
\end{frame}

% ── WARM-UP ──────────────────────────────────────────────────────────────────
\begin{frame}[t, plain]
  \navyheader{Warm-up}
  \sectionlabel{5 minutes $\cdot$ on your own}
  \begin{enumerate}\setlength{\itemsep}{8pt}
    \item Teens who eat breakfast have higher GPAs. Observational study or experiment? Does it
          show \emph{cause}?
    \item A review game in 1st period, none in 6th; 1st period scores higher. Name one
          \emph{other} reason.
    \item Use 20 slips of paper to split 20 students into two groups of 10 \emph{by chance}.
  \end{enumerate}
  \vspace{6pt}
  \begin{block}{Keep item 2 in mind}
    ``Morning'' and ``stronger class'' have a name today: \redink{confounding variables}.
  \end{block}
\end{frame}

% ── I DO — direct instruction begins ─────────────────────────────────────────
{
\setbeamercolor{background canvas}{bg=navy}
\begin{frame}[plain]
  \vspace{1.8cm}\hspace{0.8cm}%
  \begin{minipage}{0.86\textwidth}
    {\color{goldacc}\bfseries\small GUIDED NOTES \ \textbullet\ \ I DO}\\[0.7cm]
    {\color{white}\bfseries\Large Open to your Guided Notes. Fill each blank as we name it
    --- not before.}\\[0.9cm]
    {\color{skymid}\small Four terms go in the vocabulary box today: confounding variable,
    control group, placebo, single- and double-blind.}
  \end{minipage}
\end{frame}
}

% ── ROW 1: EXPERIMENT ────────────────────────────────────────────────────────
\begin{frame}[t, plain]
  \navyheader{What makes it an experiment}
  \sectionlabel{notes, Experiment $\cdot$ problems 1--4}
  \vspace{2pt}
  \begin{columns}[t]
    \begin{column}{0.48\textwidth}
      \begin{block}{The parts}
        \textbf{Experimental units} are \redink{assigned} the treatments.
        A \textbf{factor}'s levels are set on purpose; each level is a \textbf{treatment}.
        The \textbf{response} is measured \redink{after}.
      \end{block}
      \vspace{3pt}
      {\small 60 volunteers: mint gum, fruit gum, or no gum --- chosen by computer --- then
      count the words recalled.}
    \end{column}
    \begin{column}{0.48\textwidth}
      {\small\textbf{The four principles}}
      \begin{enumerate}\setlength{\itemsep}{3pt}\small
        \item \textbf{Comparison} --- at least \redink{two} groups (a control group counts).
        \item \textbf{Random assignment} --- \redink{chance} decides who gets what, so
              confounding variables balance out.
        \item \textbf{Replication} --- more than \redink{one} unit per group.
        \item \textbf{Control} --- everything else the \redink{same}.
      \end{enumerate}
    \end{column}
  \end{columns}
\end{frame}

% ── ROW 2: BLOCK IT? ─────────────────────────────────────────────────────────
\begin{frame}[t, plain]
  \navyheader{Block it?}
  \sectionlabel{notes, Block It? $\cdot$ problems 5--7}
  \vspace{2pt}
  \begin{columns}[t]
    \begin{column}{0.48\textwidth}
      \begin{block}{Completely randomized design}
        Treatments assigned to all units completely \redink{at random}.
      \end{block}
      \vspace{3pt}
      \begin{block}{Randomized block design}
        Sort units into blocks alike on a variable that affects the \redink{response}; then
        randomize \redink{within} each block.
      \end{block}
    \end{column}
    \begin{column}{0.48\textwidth}
      {\small Sports drink vs.\ water on mile times, 40 runners:}
      \begin{itemize}\setlength{\itemsep}{3pt}\small
        \item Varsity (20) $\to$ random: 10 drink, 10 water
        \item Gym class (20) $\to$ random: 10 drink, 10 water
        \item Compare mile times \emph{within} each block.
      \end{itemize}
      \vspace{6pt}
      {\small ``Block on birth month instead.'' Useful? Does birth month affect a mile time?}
    \end{column}
  \end{columns}
\end{frame}

% ── ROW 3: MATCHED PAIRS — THE CRUX ──────────────────────────────────────────
\begin{frame}[t, plain]
  \navyheader{Matched pairs: a block of two}
  \sectionlabel{notes, Matched Pairs $\cdot$ problems 8--10 $\cdot$ the whole point of today}
  \vspace{4pt}
  \begin{columns}[t]
    \begin{column}{0.46\textwidth}
      \begin{block}{Two ways to pair}
        Pair alike units; \redink{chance} picks who gets which treatment.\\
        Or each unit gets \redink{both} treatments, in a random \redink{order}.
      \end{block}
    \end{column}
    \begin{column}{0.46\textwidth}
      \begin{block}{The shoe sole}
        15 runners, one new sole and one old sole each. A coin flip picks the foot.
      \end{block}
    \end{column}
  \end{columns}
  \vspace{8pt}
  ``Each runner wears both soles, so the pairing controls everything. Skip the coin --- new
  sole on every right foot.'' \emph{Agree?}
  \vspace{8pt}
  \begin{block}{Say it, write it in the box}
    Blocking and pairing make the comparison fairer. They \redink{never replace} the coin flip.
  \end{block}
\end{frame}

% ── WE DO ────────────────────────────────────────────────────────────────────
\begin{frame}[t, plain]
  \navyheader{Guided Practice: The Zinc Lozenge}
  \sectionlabel[goldacc]{We do $\cdot$ you hold the pen $\cdot$ problems 11--14}
  120 adults with a new cold; 40 are 60 or older, and older adults have longer colds. The
  control group gets a look-alike lozenge with no zinc.
  \vspace{4pt}
  \begin{center}
  {\small\renewcommand{\arraystretch}{1.3}
  \begin{tabular}{@{} l l l @{}}
    \textbf{Block} & \textbf{Random assignment} & \textbf{Response} \\ \hline
    60 or older (40) & 20 zinc, 20 placebo & days of cold \\
    Under 60 (80) & 40 zinc, 40 placebo & days of cold \\
  \end{tabular}}
  \end{center}
  \begin{itemize}\setlength{\itemsep}{3pt}
    \item Units, treatments, response? Why block on age?
    \item Why a placebo instead of nothing? How could it be double-blind?
  \end{itemize}
\end{frame}

% ── DEBRIEF ──────────────────────────────────────────────────────────────────
\begin{frame}[t, plain]
  \navyheader{Debrief: block what matters, then flip the coin}
  \sectionlabel{8 minutes $\cdot$ pens down, papers up --- we say these aloud}
  \vspace{6pt}
  \begin{columns}[t]
    \begin{column}{0.46\textwidth}
      \begin{block}{Block on it}
        Runner experience, a person's age, the runner's own two feet: each \redink{affects the
        response}.
      \end{block}
    \end{column}
    \begin{column}{0.46\textwidth}
      \begin{block}{Then randomize}
        Who drinks, who gets zinc, which foot gets the new sole: \redink{chance decides},
        inside every block.
      \end{block}
    \end{column}
  \end{columns}
  \vspace{12pt}
  \begin{block}{Two things to say out loud}
    Why is ``the pairs are matched, so we can skip the coin'' wrong? And why does birth month
    make a useless block?
  \end{block}
\end{frame}

% ── CLOSE & START THE HOMEWORK ──────────────────────────────────────────────
\begin{frame}[t, plain]
  \navyheader{Homework --- start it now}
  \sectionlabel[goldacc]{5 items $\cdot$ scored 1 to 5 $\cdot$ due the first class after two study halls}
  \begin{itemize}\setlength{\itemsep}{5pt}
    \item Part A: two corn varieties on wet and dry fields --- the parts, why block, and a
          farmhand who wants to skip the randomizing.
    \item Part B: one multiple-choice item on a new pillow, with a sentence of justification,
          and a joint chew for older dogs: placebo and blinding.
  \end{itemize}
  \begin{block}{Silent and alone --- I am circulating}
    \emph{Part A, item 3} is the one I am reading over your shoulder. The AP Practice page is
    \textbf{extra credit, scored 1 to 5} --- work it for the reps.
  \end{block}
\end{frame}

% ── CLOSE ────────────────────────────────────────────────────────────────────
{
\setbeamercolor{background canvas}{bg=navy}
\begin{frame}[plain]
  \vspace{1.5cm}\hspace{0.8cm}%
  \begin{minipage}{0.86\textwidth}
    {\color{goldacc}\bfseries\small BEFORE YOU GO}\\[0.7cm]
    {\color{white}\bfseries\Large Block on what affects the response. Then let chance assign
    the treatments --- inside every block, every pair.}\\[0.9cm]
    {\color{skymid}\small Next: Lesson 3.5 --- Inference and Experiments: when is a difference
    between treatment groups too big to be chance alone?}
  \end{minipage}
\end{frame}
}

\end{document}
